\section{Boundary factors of torus cuts}%
\label{sec:peps_torus_cut_boundary_factors}

For a region whose lift is simply connected, the torus vectors of all
commuting closures factor across the boundary of the region through one
common boundary factor, up to a single unitary acting on the complement.

\begin{lemma}[Boundary routes through two rooted trees]%
    \label{thm:peps_regular_boundary_route}
    \lean{TNLean.PEPS.regularWalkHolonomy_inducedWalk,
        TNLean.PEPS.regularRegionTreeWalk,
        TNLean.PEPS.regularRegionTreeWalk_holonomy,
        TNLean.PEPS.regularBoundaryRoute,
        TNLean.PEPS.regularBoundaryRoute_holonomy,
        TNLean.PEPS.regularBoundaryTreeRoute,
        TNLean.PEPS.regularBoundaryTreeRoute_holonomy,
        TNLean.PEPS.regularCombinedBoundaryTransport,
        TNLean.PEPS.regularCombinedBoundaryTransport_apply,
        TNLean.PEPS.regularCombinedBoundaryTransport_tail,
        TNLean.PEPS.regularCombinedBoundaryTransport_head}
    \leanok
    \uses{thm:peps_regular_walk_holonomy,
        thm:peps_regular_region_cycle_holonomy,
        thm:peps_regular_region_tree_gauge,
        def:peps_regionComplementBoundaryConfig}
    Choose rooted spanning trees in a region $R$ and its complement $S$,
    with normalized gauges $k_R,k_S$. Let $F_R,F_S$ be their
    crossing-label maps. After rereading each crossing bond as the
    same native edge on either side, define $M=F_S\circ F_R^{-1}$.
    If $f$ joins $v_f\in R$ to $w_f\in S$, write $D_f$ for its
    directed bond operator from $v_f$ to $w_f$. Then
    \begin{align}
        M(\theta)_f & =m_f\theta_f,
                    & m_f           & =k_S(w_f)^{-1}D_fk_R(v_f).
        \label{eq:peps_boundary_route_multiplier}
    \end{align}
    Thus $D_f=u_f$ when the region endpoint is the smaller endpoint,
    and $D_f=u_f^{-1}$ when it is the larger endpoint.

    Fix a reference crossing $f_0$. Starting at the complementary root,
    construct the closed walk $p_f$ by following its tree to
    $w_{f_0}$, crossing backwards to $v_{f_0}$, following the region tree
    to $v_f$, crossing to $w_f$, and following the complementary tree
    back to its root. Its holonomy is
    \begin{align}
        H(u,p_f) & =m_fm_{f_0}^{-1}.
    \end{align}
    The construction and identity permit arbitrary bond operators.
    Inclusion of an induced-region walk in the ambient graph preserves
    its holonomy. This gives the routes for the cycle argument in
    \cite[Theorem~6.9]{Schuch2010PEPS}; the route passes through $R$.
    For commuting torus closures, the complementary loops of
    Lemma~\ref{thm:peps_torus_simply_connected_boundary_loops} have
    the same winding, and hence the required holonomy.
\end{lemma}

\begin{proof}\leanok
    The two boundary equivalences multiply at the region and
    complementary endpoints, giving the crossing multiplier
    (\ref{eq:peps_boundary_route_multiplier}).
    Along each chosen tree path from $v$ to $w$, the original bond
    holonomy is $k(w)k(v)^{-1}$. Multiplying the five successive
    portions of $p_f$ in reverse traversal order and cancelling the
    root gauges gives $m_fm_{f_0}^{-1}$. Induced-walk inclusion retains
    each ordered native edge and hence each directed factor.
\end{proof}

\begin{lemma}[Region walk crossing numbers from geometric lifts]%
    \label{thm:peps_region_walk_winding}
    \lean{TNLean.PEPS.torusWalkWinding_of_isTorusRegionIntegerLift,
        TNLean.PEPS.torusWalkWinding_eq_of_isSimplyConnected}
    \leanok
    \uses{thm:peps_torus_walk_winding,
        thm:peps_torus_region_integer_lift_gauge,
        thm:peps_regular_boundary_route,
        thm:peps_torus_region_realization}
    Let $L:R\to\mathbb Z^2$ be an integer lift preserving each native
    horizontal and vertical unit step, and write the torus periods as
    $W$ and $H$. For any induced-region walk $p:v\to t$,
    its ambient seam crossing numbers are
    \begin{align}
        N(p)=\left(
        \left\lfloor\frac{L_1(t)}{W}\right\rfloor
        -\left\lfloor\frac{L_1(v)}{W}\right\rfloor,
        \left\lfloor\frac{L_2(t)}{H}\right\rfloor
        -\left\lfloor\frac{L_2(v)}{H}\right\rfloor
        \right).
        \label{eq:peps_region_deck_winding}
    \end{align}
    As auxiliary interior-walk identities for the path argument
    in \cite[proof of Theorem~6.9]{Schuch2010PEPS}, if the closed-cell realization of $R$ is simply connected,
    any two induced-region walks with the same endpoints have equal
    seam crossing numbers; no integer lift is assumed in this case.
\end{lemma}

\begin{proof}\leanok
    In the additive group $\mathbb Z^2$, written
    multiplicatively with the universal commuting closure elements
    $g=(0,-1)$ and $h=(1,0)$, the walk holonomy is $N(p)$ by
    Lemma~\ref{thm:peps_torus_walk_winding}. The lift gauge is the pair
    of integer deck coordinates. Its native gradient is the
    closure assignment on every internal edge, so its walk holonomy
    is the difference of the endpoint deck coordinates. Inclusion of
    the induced walk preserves the native holonomy, giving
    (\ref{eq:peps_region_deck_winding}). In the simply connected case,
    Lemma~\ref{thm:peps_torus_region_realization} gives the lift;
    applying (\ref{eq:peps_region_deck_winding}) to both walks gives
    equality.
\end{proof}

\begin{lemma}[Complement cycle words from torus winding data]%
    \label{thm:peps_torus_complement_cycle_words}
    \lean{TNLean.PEPS.regularBoundaryCrossingAdj,
        TNLean.PEPS.regularBoundaryInsideVertex,
        TNLean.PEPS.regularBoundaryOutsideVertex,
        TNLean.PEPS.regularCombinedBoundaryTransport_insideOutside,
        TNLean.PEPS.torusBoundaryTreeRoute,
        TNLean.PEPS.torusBoundaryTreeRoute_holonomy,
        TNLean.PEPS.regularRegionCycleWord_eval_torusClosure_of_winding_eq}
    \leanok
    \uses{thm:peps_regular_boundary_route,
        thm:peps_regular_region_cycle_holonomy,
        thm:peps_torus_walk_winding}
    On a native torus of width and height at least three, choose rooted
    spanning trees in $R$ and $S=R^c$. For each native crossing bond $f$,
    membership determines its inside endpoint $v_f$, its outside endpoint
    $w_f$, and its directed step from $v_f$ to $w_f$. Choose a reference
    bond $f_0$, and form the two-tree route $q_f$ of
    Lemma~\ref{thm:peps_regular_boundary_route} from these endpoints.
    For arbitrary native bond operators its holonomy is
    \begin{align}
        H(u,q_f) & =m_fm_{f_0}^{-1},
                 & m_f               & =(F_S\circ F_R^{-1})(\mathbf 1)_f.
    \end{align}

    Suppose loops $p_f$ in the induced complementary graph,
    based at its normalized root, have been chosen independently of the
    closure labels. Assume their ambient seam crossing numbers satisfy
    \begin{align}
        \operatorname{wind}(p_f) & =\operatorname{wind}(q_f).
    \end{align}
    For every commuting pair $g,h$, let $u^{g,h}$ be the native
    closure assignment, $\omega_S^{g,h}$ its complementary tree residuals,
    and $m_f^{g,h}$ its combined crossing multipliers. Then the genuine
    complementary walk words satisfy
    \begin{align}
        \operatorname{ev}_{\omega_S^{g,h}}W_{T_S}(p_f)
         & =m_f^{g,h}(m_{f_0}^{g,h})^{-1}.
    \end{align}
    Thus the same words give the relative identity simultaneously for
    all commuting closure pairs. This is a geometry-to-group identity
    for the accessible complement argument in
    \cite[Theorem~6.9]{Schuch2010PEPS}, local source lines 1935--1990.
    It neither assumes a group-valued relative holonomy equation nor
    asserts existence of the complementary loops with these winding numbers;
    the remaining scope is recorded in
    \cite{gap:rmp_peps_quantum_double_g_isometry}.
\end{lemma}

\begin{proof}\leanok
    Region membership selects the two endpoints of every crossing bond.
    The directed factor is its original operator or its inverse according
    to that selection, so the two-tree route has the stated
    combined multiplier ratio. Evaluation of a complementary word at
    its tree residual equals the original root-loop holonomy.
    Inclusion in the ambient graph preserves that holonomy. Equal winding
    numbers give equal commuting-closure holonomies, identifying the
    complementary loop holonomy with the reference-route ratio.
\end{proof}

\begin{lemma}[Cycle-controlled permutation of boundary coordinates]%
    \label{thm:peps_cycle_controlled_boundary}
    \lean{TNLean.PEPS.regularCycleControlledBoundaryMatrix_mem_unitaryGroup,
        TNLean.PEPS.regularCycleControlledBoundary,
        TNLean.PEPS.regularCycleControlledBoundaryMatrix,
        TNLean.PEPS.regularCycleControlledBoundaryMatrix_mul_coordinates,
        TNLean.PEPS.regularComplementWalkControlMatrix,
        TNLean.PEPS.regularComplementWalkControlMatrix_mem_unitaryGroup}
    \leanok
    \uses{def:peps_regular_region_coordinates}
    Fix spanning-tree physical coordinates $(y,a,r,z)$ of a finite
    region and any function $W:G^C\to G^{\partial R}$ of its cycle
    labels. The permutation
    $\sigma_W(y,a,r,z)=(W(z)^{-1}y,a,r,z)$ is reversible, with inverse
    $(y,a,r,z)\mapsto(W(z)y,a,r,z)$, and its matrix $U_W$ on the
    original half-edge physical basis is unitary. In particular, fixed
    complementary walks define such a matrix by evaluating their
    cycle words; this construction is independent of all bond assignments.

    For the canonical block $M_R^u$, let $b_u$ be its boundary
    map and $\omega$ its cycle residual tuple. Its transformed
    coefficients satisfy
    \begin{align}
        (U_WM_R^u)\bigl(\alpha(y,a,r,z),\theta\bigr)
         & =|G|^{-|R|}\sum_{x\in G}
        \mathbf 1_{\{y=W(z)^{-1}x b_u(\theta)\}}
        \mathbf 1_{\{z=x\omega x^{-1}\}}.
        \label{eq:peps_cycle_controlled_coefficients}
    \end{align}
    In this algebraic coordinate operation associated with
    \cite[\S6.3]{Schuch2010PEPS}, the cycle constraint is retained without any assumption on the
    relation between $W$ and the crossing maps.
\end{lemma}
\begin{proof}\leanok
    \uses{thm:peps_regular_twisted_projector_coordinates}
    The inverse translates only $y$, so it leaves its own controlling
    tuple $z$ fixed. Group cancellation proves the inverse identities.
    A permutation of the physical group basis is unitary. Its action
    on coefficient functions is precomposition with $\sigma_W^{-1}$;
    substituting $W(z)y$ for $y$ in
    (\ref{eq:peps_twisted_projector_coordinates}) gives
    (\ref{eq:peps_cycle_controlled_coefficients}).
\end{proof}

\begin{lemma}[Combined boundary shift removal from given complement walks]%
    \label{thm:peps_cycle_controlled_relative_walks}
    \lean{TNLean.PEPS.regularCycleControlledBoundaryMatrix_mul_coordinates_of_combined_relative_walk_words}
    \leanok
    \uses{thm:peps_cycle_controlled_boundary,
        thm:peps_regular_region_cycle_holonomy,
        thm:peps_regular_boundary_route}
    Let $S$ be the complement of $R$, and identify their crossing bonds
    by the same native edges. Choose a region gauge $k_R$ and a rooted
    spanning tree in $S$, with normalized gauge $k_S$ and cycle residuals
    $\omega_S$. Let $F_R,F_S$ be the boundary maps and set
    $M=F_S\circ F_R^{-1}$ under this crossing identification. Write
    $m_f=M(\mathbf 1)_f$, so that $M(\theta)_f=m_f\theta_f$.
    Fix a reference crossing bond $f_0$ and a walk $p_f$ entirely
    in the induced graph on $S$ for each crossing bond $f$. Assume
    \begin{align}
        \operatorname{ev}_{\omega_S}W_T(p_f) & =m_fm_{f_0}^{-1}.
        \label{eq:peps_control_relative_walks}
    \end{align}
    Define $W(z)_f=\operatorname{ev}_zW_T(p_f)$. With the complementary
    boundary input reread from $F_R^{-1}\theta$, the coefficients
    satisfy
    \begin{align}
        (U_WM_S^u)\bigl(\alpha_S(y,a,r,z),F_R^{-1}\theta\bigr)
         & =|G|^{-|S|}\sum_{x\in G}
        \mathbf 1_{\{y_f=xm_{f_0}\theta_f\ (f\in\partial S)\}}
        \mathbf 1_{\{z=x\omega_S x^{-1}\}}.
        \label{eq:peps_control_common_transport}
    \end{align}
    Here $\theta_f$ denotes the label on the corresponding crossing bond
    of $R$. This is an auxiliary conditional statement in the complement
    argument of \cite[\S6.3]{Schuch2010PEPS}. Constructing the required
    complement walks with (\ref{eq:peps_control_relative_walks}) remains
    the geometric step described in
    \cite{gap:rmp_peps_quantum_double_g_isometry}.
\end{lemma}
\begin{proof}\leanok
    At the reindexed input, the complementary boundary map
    is $M(\theta)_f=m_f\theta_f$. On the cycle support, conjugation
    of the constructed words gives
    \begin{align}
        W(x\omega_S x^{-1})_f   & =xm_fm_{f_0}^{-1}x^{-1}, \\
        W(z)_f^{-1}xm_f\theta_f & =xm_{f_0}\theta_f.
        \label{eq:peps_control_relative_cancellation}
    \end{align}
    Substituting~(\ref{eq:peps_control_relative_cancellation}) into
    (\ref{eq:peps_cycle_controlled_coefficients}) retains the cycle indicator.
\end{proof}

\begin{lemma}[Relative boundary factor of controlled complementary blocks]%
    \label{thm:peps_controlled_complement_boundary_factor}
    \lean{TNLean.PEPS.regularControlledComplementMatrix_factorization,
        TNLean.PEPS.regularControlledComplementMatrix,
        TNLean.PEPS.regularControlledComplementMatrix_sum_factorization}
    \leanok
    \uses{thm:peps_regular_boundary_invariants}
    Assume the complement walks satisfy the relative-word
    identities of (\ref{eq:peps_control_relative_walks}). Number the
    $b=n+1$ crossing bonds and use reference boundary coordinates
    $y=(t,d)$ and $\theta=(a,r)$, where $d,r\in G^n$. For the
    controlled complementary matrix at the reindexed input
    $F_R^{-1}\theta$, put $x=t a^{-1}m_{f_0}^{-1}$. Its coefficient is
    \begin{align}
        |G|^{-|S|}\mathbf 1_{\{d=r\}}
        \mathbf 1_{\{z=x\omega_Sx^{-1}\}}.
    \end{align}
    The normalized boundary factor is
    $|G|^{-n/2}\mathbf 1_{\{d=r\}}$.
    Its remaining factor is
    $|G|^{n/2-|S|}\mathbf 1_{\{z=x\omega_Sx^{-1}\}}$.
    The trees, numbering, and control walks may be fixed before an
    arbitrary finite family of bond assignments. The same factor then
    separates from every coherent linear combination of their
    complementary matrices; the residual coefficients are combined
    with the same weights.

    This statement retains the virtual reference label $a$ in the
    residual coefficient. It is an auxiliary coefficient statement,
    rather than a conclusion about the full physical bipartite state;
    for commuting torus closures, the complementary loops of
    Lemma~\ref{thm:peps_torus_simply_connected_boundary_loops} realize
    its geometric walk condition.
\end{lemma}
\begin{proof}\leanok
    \uses{thm:peps_cycle_controlled_relative_walks}
    Equality under a common left translation is equivalent to equality
    of all relative labels and equality of the reference labels.
    The latter fixes $x=t a^{-1}m_{f_0}^{-1}$ uniquely. The finite
    translation sum therefore has one surviving term. Multiplication
    by $|G|^{n/2}$ isolates the normalized maximally entangled
    coefficient, and linearity gives the coherent sum identity.
\end{proof}

\begin{lemma}[Common boundary factor of controlled canonical cuts]%
    \label{thm:peps_controlled_projector_cut_boundary_factor}
    \lean{TNLean.PEPS.regularControlledProjectorCutMatrix_eq_native_contraction,
        TNLean.PEPS.regularControlledProjectorCutMatrix,
        TNLean.PEPS.regularControlledProjectorCutMatrix_coordinates,
        TNLean.PEPS.regularReferenceRegionAncilla,
        TNLean.PEPS.regularReferenceRegionAncilla_dotProduct,
        TNLean.PEPS.regularControlledProjectorCutMatrix_sum_factorization}
    \leanok
    \uses{thm:peps_controlled_complement_boundary_factor, def:peps_regular_region_coordinates}
    Retain the same complement walk identities, and suppose
    every cycle residual in $R$ is the identity. Contract the
    two canonical region matrices, applying the fixed controlled
    physical permutation only in $S$. Reindexing their common boundary
    sum by $F_R$ preserves the original contraction.

    Let $d_R,d_S\in G^n$ be the relative boundary labels. In the
    spanning-tree physical coordinates the contracted
    coefficient is
    \begin{align}
        |G|^{-|R|-|S|}\mathbf 1_{\{d_R=d_S\}}
        \mathbf 1_{\{z_R=\mathbf 1\}}
        \sum_{x\in G}\mathbf 1_{\{z_S=x\omega_Sx^{-1}\}}.
        \label{eq:peps_controlled_canonical_cut_factor}
    \end{align}
    In particular neither reference boundary label nor the common
    multiplier $m_{f_0}$ occurs in the result.

    Let $N_R$ be the number of choices of the reference boundary label,
    internal reference labels, and normalized vertex labels in $R$.
    The fixed region ancillary vector is
    $\phi_R=N_R^{-1/2}\mathbf 1_{\{z_R=\mathbf 1\}}$; it has norm
    one. For any finite sector family obeying the stated
    residual and walk conditions, and any coefficients $\lambda_i$,
    the coherent cut sum has the exact factorization
    \begin{align}
        \Omega_{|G|^n}(d_R,d_S)\,\phi_R\,
        \biggl(|G|^{n/2}\sqrt{N_R}\,|G|^{-|R|-|S|}
        \sum_i\lambda_i\sum_{x\in G}
        \mathbf 1_{\{z_S=x\omega_{S,i}x^{-1}\}}\biggr).
    \end{align}
    Both normalized factors are fixed before the sector family and
    its coefficients; all sector dependence is on the complementary
    cycle coordinates. These are auxiliary canonical-cut statements
    associated with \cite[Theorem~6.9]{Schuch2010PEPS}. The
    geometric conditions and transfer to the original physical
    tensors remain separate steps, as recorded in
    \cite{gap:rmp_peps_quantum_double_g_isometry}.
\end{lemma}
\begin{proof}\leanok
    \uses{thm:peps_regular_twisted_projector_coordinates}
    Changing the common boundary summation variable by its
    reversible region boundary map and substituting the two
    spanning-tree coefficient formulas fixes the
    shared boundary configuration for each complementary translation via the complementary boundary indicator. Its relative labels equal $d_S$,
    and the interior translation sum gives
    $\mathbf 1_{\{d_R=d_S\}}$. This proves
    (\ref{eq:peps_controlled_canonical_cut_factor}). There are exactly
    $N_R$ unrestricted region ancillary configurations, while its
    cycle tuple is fixed, proving normalization. Extracting the two
    square-root normalization factors and using linearity yields the
    finite sector sum.
\end{proof}

\begin{lemma}[A common canonical torus cut factor from simple connectedness and winding]%
    \label{thm:peps_torus_controlled_boundary_factor}
    \lean{TNLean.PEPS.exists_unitary_torusControlledBoundaryFactor_of_isSimplyConnected_and_winding}
    \leanok
    \uses{thm:peps_controlled_projector_cut_boundary_factor,
        thm:peps_torus_region_realization,
        thm:peps_torus_region_integer_lift_gauge,
        thm:peps_torus_complement_cycle_words}
    Let both torus dimensions be at least three. Suppose the
    closed-cell realization of $R$ is simply connected. Choose rooted
    spanning trees in $R$ and $S$, number their nonempty crossing bonds
    by $\{0,\ldots,n\}$, and choose fixed complementary root
    loops $p_f$. Assume their ambient torus windings equal those of
    the native two-tree routes from the reference crossing to $f$.

    There exists one unitary $U$ on the original complementary
    half-edge physical basis, chosen before all closure sectors and
    coefficients, with the following property. For every finite family
    of commuting pairs $(g_i,h_i)$ and coefficients $\lambda_i$, form
    the canonical cut contractions
    \begin{align}
        K_i(\alpha,\beta)
         & =\sum_{\theta}
        M_R^{g_i,h_i}(\alpha,\theta)
        (UM_S^{g_i,h_i})(\beta,\theta),
        \label{eq:peps_torus_canonical_cut_contraction}
    \end{align}
    where the same native crossing edges are reread on the two sides.
    In the fixed spanning-tree and relative boundary coordinates,
    their coherent sum is
    \begin{align}
        \sum_i\lambda_i K_i
         & =\Omega_{|G|^n}\otimes\phi_R\otimes\chi_{\lambda},
    \end{align}
    with the same normalized maximally entangled factor and the same
    normalized region ancillary vector as in
    Lemma~\ref{thm:peps_controlled_projector_cut_boundary_factor}.
    The remaining complementary coefficient is
    \begin{align}
        \chi_{\lambda}(z_S)
         & =|G|^{n/2}\sqrt{N_R}\,|G|^{-|R|-|S|}
        \sum_i\lambda_i\sum_{x\in G}
        \mathbf 1_{\{z_S=x\omega_{S,i}x^{-1}\}}.
    \end{align}
    Thus no group-valued flatness, relative-word identity, mixed Gram
    formula, or tensor factorization is assumed as a hypothesis.
    This is an auxiliary canonical-cut consequence in the argument of
    \cite[Theorem~6.9]{Schuch2010PEPS}. Existence of the required
    complementary loops for the paper's disk geometry and transfer
    to the original physical site tensors remain separate steps;
    see \cite{gap:rmp_peps_quantum_double_g_isometry}.
\end{lemma}
\begin{proof}\leanok
    Simple connectedness of the realization gives an integer
    coordinate lift. For every commuting pair its normalized gauge
    gives identity region cycle residuals. Equality of the given
    loop windings with the native routes gives the complement
    relative-word identities simultaneously for all commuting pairs.
    Evaluating the fixed walk words constructs $U$, and applying the
    coherent canonical-cut factorization before reindexing the shared
    boundary sum recovers the original contraction
    (\ref{eq:peps_torus_canonical_cut_contraction}).
\end{proof}
